\chapter{Family Planning and Contraception}
\index{Family Planning and Contraception}

\section*{Definition and Overview}
Family planning is the process of deciding if, when, and how many children to have, and using the methods available to achieve that goal. Contraception, the deliberate prevention of pregnancy, is a central part of family planning and encompasses a wide range of methods: hormonal options such as the pill, implant, and intrauterine devices; barrier methods such as condoms and diaphragms; permanent methods such as vasectomy and tubal ligation; and fertility awareness approaches. Modern contraception also protects against sexually transmitted infections when condoms are used, and allows women and couples to space pregnancies for health, social, and economic reasons. This chapter outlines the available methods, how to choose between them, and their effectiveness, benefits, and risks.

\section*{Epidemiology}
Contraceptive use is widespread in many countries and most of the world. Around two-thirds of women of reproductive age use some form of contraception, with the oral contraceptive pill, condoms, and intrauterine devices among the most common choices. Rates of unintended pregnancy remain significant, particularly among younger women and those using less reliable methods inconsistently, and about a quarter of local pregnancies are unplanned. Use of long-acting reversible contraception (LARC), such as implants and intrauterine devices, has increased in recent decades because of its high effectiveness and convenience, although uptake still lags behind many comparable countries. Access to affordable, effective contraception is recognised internationally as a key determinant of maternal and child health outcomes.

\section*{Aetiology and Pathophysiology}
Contraceptive methods prevent pregnancy through several distinct mechanisms.
\begin{itemize}
    \item \textbf{Hormonal contraceptives:} oestrogen and progestogen combinations suppress ovulation, thicken cervical mucus to impede sperm, and alter the endometrial lining; progestogen-only methods act mainly through cervical mucus and endometrial changes, with variable ovulation suppression
    \item \textbf{Intrauterine devices:} the copper IUD is toxic to sperm and eggs and impairs their transport, while the levonorgestrel IUD thins the endometrium and thickens cervical mucus, and suppresses ovulation in some women
    \item \textbf{Barrier methods:} male and female condoms and diaphragms physically prevent sperm from reaching the egg; condoms additionally prevent the transmission of most sexually transmitted infections
    \item \textbf{Permanent methods:} vasectomy blocks the vas deferens so sperm cannot enter the ejaculate, while tubal ligation or removal blocks or removes the fallopian tubes so the egg cannot meet sperm
    \item \textbf{Fertility awareness:} relies on identifying the fertile window through menstrual cycle tracking of temperature, cervical mucus, and calendar calculations, then avoiding intercourse or using barriers during that window
\end{itemize}

\section*{Clinical Features}
Contraception is used by healthy people and does not produce a set of symptoms. The relevant features are the method's characteristics and acceptability to the user.
\begin{itemize}
    \item The effectiveness of each method, usually expressed as the failure rate with typical use versus perfect use
    \item The convenience and reversibility of the method, including whether it requires daily attention or a long-acting device
    \item Side effects, which for hormonal methods may include irregular bleeding, breast tenderness, mood changes, and headaches
    \item Non-contraceptive benefits, such as lighter periods, reduced period pain, and improved acne with some hormonal methods
    \item The level of protection against sexually transmitted infections, which only condoms provide reliably
    \item The user's personal health history, preferences, and future pregnancy plans
\end{itemize}

\section*{Differential Diagnosis}
When a person presents to discuss contraception or with a concern about it, the main considerations are:
\begin{itemize}
    \item \textbf{Pregnancy:} always exclude pregnancy before starting or changing a method, and consider it if a user has missed a period or made a method error
    \item \textbf{Method-related side effects:} breakthrough bleeding, mood symptoms, and other complaints must be distinguished from unrelated conditions
    \item \textbf{Medical conditions affecting method choice:} migraine with aura, hypertension, venous thromboembolism history, and liver disease influence whether oestrogen-containing methods are appropriate
    \item \textbf{Sexually transmitted infections:} symptoms of discharge, pain, or sores require testing regardless of contraceptive use
    \item \textbf{Gynaecological conditions:} heavy bleeding, pelvic pain, or amenorrhoea warrant assessment independent of contraceptive use
\end{itemize}

\section*{When to Seek Medical Care}
People should seek advice when choosing a contraceptive method, especially if they have medical conditions, are over 35 and smoke, or have a family history of blood clots or breast cancer. A doctor or family planning clinic can help match a method to individual needs and health circumstances.

\textbf{Contact your local emergency services for:}
\begin{itemize}
    \item Severe abdominal or pelvic pain after IUD insertion, which may indicate perforation or infection
    \item Sudden severe chest pain, breathlessness, coughing up blood, or painful swollen calf, which may indicate a blood clot in hormonal contraceptive users
    \item Sudden severe headache, visual disturbance, or weakness, which may indicate stroke
    \item Signs of anaphylaxis after device insertion or a new medicine
\end{itemize}

\section*{Diagnosis and Investigations}
Most contraceptive consultations require no investigations, but assessment is individualised.
\begin{itemize}
    \item \textbf{History:} menstrual pattern, pregnancy plans, medical conditions, medications, family history, and prior contraceptive experiences
    \item \textbf{Pregnancy test:} to exclude pregnancy before initiating a method when there is any doubt
    \item \textbf{Blood pressure measurement:} required before prescribing oestrogen-containing methods
    \item \textbf{STI screening:} recommended, particularly for those at risk, as part of comprehensive sexual health care
    \item \textbf{Pelvic examination:} sometimes needed for diaphragm fitting or IUD insertion, but not required for the pill or implant
\end{itemize}

\section*{Management}
\begin{itemize}
    \item \textbf{Combined hormonal methods:} pill, vaginal ring, and patch contain oestrogen and progestogen; highly effective with perfect use, but require daily, weekly, or monthly attention
    \item \textbf{Progestogen-only methods:} the mini-pill, injection, and implant suit those who cannot take oestrogen; the implant is among the most effective methods available
    \item \textbf{Intrauterine devices:} copper IUDs provide non-hormonal contraception for up to 10 years; levonorgestrel IUDs last 5--8 years and often reduce bleeding
    \item \textbf{Barrier methods:} male and female condoms prevent pregnancy and most STIs; diaphragms and cervical caps require fitting and spermicide
    \item \textbf{Permanent contraception:} vasectomy is a minor day procedure with a low failure rate; tubal ligation is performed under general anaesthesia
    \item \textbf{Emergency contraception:} the copper IUD or oral emergency pills can prevent pregnancy when used within set timeframes after unprotected sex
    \item \textbf{Fertility awareness:} effective for well-motivated couples with regular cycles, but demands strict daily adherence
\end{itemize}

\section*{Complications}
\begin{itemize}
    \item Unintended pregnancy from method failure or inconsistent use, particularly with user-dependent methods
    \item Hormonal side effects including irregular bleeding, breast tenderness, mood changes, weight gain, and reduced libido
    \item Small increased risk of venous thromboembolism and stroke with combined hormonal methods, higher in smokers and older women
    \item Uterine perforation or pelvic infection, rarely, at IUD insertion
    \item Method failure, including ectopic pregnancy, which must be considered with any pregnancy while using an IUD or progestogen-only method
    \item Risks of anaesthesia and surgery with permanent methods
\end{itemize}

\section*{Prognosis and Outcomes}
All commonly used contraceptive methods are safe and effective when used correctly, and serious complications are rare. Fertility returns rapidly after stopping most methods: immediately after the pill, injection, implant, or barrier methods, and within days or weeks of IUD removal, although there may be a short delay after injectables. Long-acting reversible contraception has the highest typical-use effectiveness because it removes the risk of user error. Choosing a method that suits a person's lifestyle, health, and preferences markedly improves satisfaction and continuation, and couples who use contraception consistently can space or prevent pregnancies reliably.

\section*{Prevention and Self-Care}
\begin{itemize}
    \item Discuss contraceptive options with a doctor, nurse, or family planning clinic before needing them, rather than in an emergency
    \item Use condoms in addition to another method to prevent sexually transmitted infections
    \item Take daily methods at the same time each day and know what to do after missed pills or late injections
    \item Keep a record of IUD or implant insertion and review dates
    \item Seek help early if side effects are troublesome rather than stopping abruptly, since a doctor can offer alternatives
    \item Know where to obtain emergency contraception and its timeframes
    \item Reassess contraception as circumstances change, including after childbirth, while breastfeeding, and approaching menopause
\end{itemize}

\section*{Sources and Further Reading}
The following reputable sources provide further information for healthcare students and patients seeking reliable, current advice about family planning and contraception.
\begin{itemize}
    \item Healthdirect Australia. \textit{Contraception: methods and how to choose.} https://www.healthdirect.gov.au/contraception
    \item Family Planning NSW. \textit{Contraception fact sheets.} https://www.fpnsw.org.au/
    \item NHS. \textit{Contraception guide.} https://www.nhs.uk/conditions/contraception/
    \item Mayo Clinic. \textit{Birth control options.} https://www.mayoclinic.org/
    \item MSD Manual. \textit{Contraception: overview for clinicians.} https://www.msdmanuals.com/
    \item World Health Organization. \textit{Family planning and contraception.} https://www.who.int/
\end{itemize}
